\section{Conclusion}\label{sec:con}
In this work, we introduced a new paradigm for human video generation inspired by the dual-system theory of human cognition. We argued that existing methods primarily simulate reactive "System 1" thinking, failing to align motion with high-level intent. We proposed OmniHuman-1.5, a framework that additionally models deliberative "System 2" processes through two key innovations: an MLLM-based agent for semantic planning and a specialized MMDiT architecture with a novel pseudo last frame strategy to fuse multimodal signals.
Experiments show our approach generates more expressive and logically consistent results, which users significantly preferred for their naturalness and plausibility. By demonstrating this framework's effectiveness, even extending it to multi-person scenarios, we believe simulating cognitive agency offers a new perspective for creating the next generation of lifelike digital humans.


\section{Broader Impact}
Our core contribution in this work is to introduce a novel paradigm for video avatar generation. By simulating a dual-system cognitive framework, our model achieves a new level of expressive capability and logical coherence in motion, moving beyond the limitations of single-process generation. While this advancement opens up exciting possibilities for creative applications like AI-driven film production and music videos, we are acutely aware of the potential for misuse associated with highly realistic avatar technologies.
To address these ethical concerns, we advocate for a robust framework of responsible deployment. Although current results may still bear subtle artifacts of AI generation, which can serve as a minor deterrent, proactive safeguards are essential. We strongly recommend the following measures: (1) applying prominent, visible watermarks to all generated content to clearly label it as AI-generated; (2) implementing filtering algorithms to reject inappropriate or malicious input prompts and to review output content; and (3) embedding traceable, invisible watermarks to ensure accountability and aid in source identification if misuse occurs. By integrating these safety protocols, we can help ensure that our technology fosters creativity while minimizing the risks of malicious applications such as fraud or disinformation.



\section{Acknowledgment}
We thank Ashley Liang, Bei Li, Xu Song, Jie Ling, Wei Han and Xi Yang for their help with the evaluation. We also thank the ByteDance Seed Video Generation Team for their valuable discussions regarding model training.
